% ==============================================================================
% annexes/glossary.tex — UE SIS589
% Annexe B : Glossaire Bilingue FR/EN
% ==============================================================================

\chapter{Glossaire Bilingue Français~/ Anglais}
\label{ann:glossary}

\begin{boxinfo}
Glossaire des termes techniques utilisés dans ce cours. La langue d'usage
en cybersécurité opérationnelle étant souvent l'anglais, les termes
anglophones sont systématiquement mentionnés. Un terme en italique indique
un emprunt à l'anglais désormais d'usage courant dans la communauté
francophone.
\end{boxinfo}

\renewcommand{\arraystretch}{1.25}

% ─── A ────────────────────────────────────────────────────────────────────────
\section*{A}

\begin{table}[H]\centering\small\rowcolors{2}{TableauPair}{TableauImpair}
\begin{tabular}{p{4.5cm}p{3.5cm}p{6.5cm}}
\toprule
\tableauHeader{Terme FR} & \tableauHeader{EN} & \tableauHeader{Définition} \\
\midrule
Admissibilité & Admissibility &
Qualité d'une preuve pouvant être présentée et acceptée par une juridiction. \\
Agent (SIEM) & Agent &
Processus léger installé sur un hôte pour collecter et transmettre les logs
au serveur SIEM. \\
Alerte & Alert &
Notification générée par le SIEM lorsqu'une règle de corrélation est satisfaite. \\
Artefact (forensique) & Artifact &
Trace ou vestige numérique laissé par une activité système ou utilisateur,
utilisable comme preuve. \\
\textit{Assume breach} & Assume breach &
Paradigme de sécurité postulant que l'attaquant est déjà présent dans le SI. \\
\bottomrule
\end{tabular}
\end{table}

% ─── B ────────────────────────────────────────────────────────────────────────
\section*{B}

\begin{table}[H]\centering\small\rowcolors{2}{TableauPair}{TableauImpair}
\begin{tabular}{p{4.5cm}p{3.5cm}p{6.5cm}}
\toprule
\tableauHeader{Terme FR} & \tableauHeader{EN} & \tableauHeader{Définition} \\
\midrule
\textit{Beacon} (C2) & Beacon &
Communication périodique entre un implant et son infrastructure C2. \\
Balisage & Timestamping &
Apposition d'un horodatage fiable sur un document ou une preuve numérique. \\
\textit{Baseline} & Baseline &
État de référence (comportement normal) servant de base à la détection
d'anomalies. \\
Bloc non alloué & Unallocated block &
Secteur disque non actuellement utilisé par le système de fichiers~;
peut contenir des données antérieures récupérables. \\
Bloqueur en écriture & Write blocker &
Dispositif matériel ou logiciel empêchant toute écriture sur un support
pendant son acquisition forensique. \\
\bottomrule
\end{tabular}
\end{table}

% ─── C ────────────────────────────────────────────────────────────────────────
\section*{C}

\begin{table}[H]\centering\small\rowcolors{2}{TableauPair}{TableauImpair}
\begin{tabular}{p{4.5cm}p{3.5cm}p{6.5cm}}
\toprule
\tableauHeader{Terme FR} & \tableauHeader{EN} & \tableauHeader{Définition} \\
\midrule
\textit{Carving} de fichiers & File carving &
Récupération de fichiers supprimés par recherche de signatures dans l'espace
non alloué, sans s'appuyer sur les métadonnées du système de fichiers. \\
Chaîne de custody & Chain of custody &
Traçabilité documentée de chaque accès à une pièce à conviction depuis
sa collecte jusqu'à sa présentation en justice. \\
Commande et Contrôle & Command and Control (C2) &
Infrastructure utilisée par l'attaquant pour communiquer avec les implants
déployés sur les systèmes compromis. \\
Confinement & Containment &
Mesures visant à limiter la propagation d'un incident sans altérer les preuves. \\
Corrélation & Correlation &
Processus SIEM d'analyse croisée d'événements hétérogènes pour identifier
des patterns d'attaque. \\
\textit{Cyber Threat Intelligence} & CTI &
Renseignement actionnabel sur les menaces, produit à partir de l'analyse
d'informations sur les adversaires. \\
\bottomrule
\end{tabular}
\end{table}

% ─── D ────────────────────────────────────────────────────────────────────────
\section*{D}

\begin{table}[H]\centering\small\rowcolors{2}{TableauPair}{TableauImpair}
\begin{tabular}{p{4.5cm}p{3.5cm}p{6.5cm}}
\toprule
\tableauHeader{Terme FR} & \tableauHeader{EN} & \tableauHeader{Définition} \\
\midrule
DKOM & Direct Kernel Object Manipulation &
Technique de rootkit modifiant directement les structures du noyau
pour dissimuler des processus. \\
\textit{Dwell time} & Dwell time &
Durée entre la compromission initiale d'un système et sa détection par
les équipes de sécurité. \\
Dump mémoire & Memory dump &
Copie bit à bit du contenu de la mémoire vive d'un système à un instant~T. \\
\bottomrule
\end{tabular}
\end{table}

% ─── E ────────────────────────────────────────────────────────────────────────
\section*{E}

\begin{table}[H]\centering\small\rowcolors{2}{TableauPair}{TableauImpair}
\begin{tabular}{p{4.5cm}p{3.5cm}p{6.5cm}}
\toprule
\tableauHeader{Terme FR} & \tableauHeader{EN} & \tableauHeader{Définition} \\
\midrule
EDR & Endpoint Detection and Response &
Agent hôte collectant la télémétrie comportementale et permettant la
réponse distante~: isolation réseau, kill de processus. \\
Éradication & Eradication &
Phase de réponse aux incidents consistant à supprimer la cause racine~:
malwares, persistances, comptes créés. \\
\textit{Event ID} (EID) & Event ID &
Identifiant numérique d'un type d'événement dans les journaux Windows. \\
Exfiltration & Exfiltration &
Transfert non autorisé de données depuis le système cible vers
une infrastructure contrôlée par l'attaquant. \\
\bottomrule
\end{tabular}
\end{table}

% ─── F ────────────────────────────────────────────────────────────────────────
\section*{F}

\begin{table}[H]\centering\small\rowcolors{2}{TableauPair}{TableauImpair}
\begin{tabular}{p{4.5cm}p{3.5cm}p{6.5cm}}
\toprule
\tableauHeader{Terme FR} & \tableauHeader{EN} & \tableauHeader{Définition} \\
\midrule
Fatigue d'alerte & Alert fatigue &
État d'épuisement des analystes SOC dû à un taux de faux positifs élevé,
conduisant à ignorer les vraies menaces. \\
FIM & File Integrity Monitoring &
Surveillance continue des modifications apportées aux fichiers systèmes
critiques~; détection d'altérations non autorisées. \\
Forensique numérique & Digital forensics &
Science d'identification, préservation, analyse et présentation des
preuves numériques de manière légalement recevable. \\
Faux positif & False positive (FP) &
Alerte signalant une menace qui n'existe pas réellement. \\
Faux négatif & False negative (FN) &
Menace réelle non détectée par les outils de sécurité. \\
\bottomrule
\end{tabular}
\end{table}

% ─── H-I ──────────────────────────────────────────────────────────────────────
\section*{H -- I}

\begin{table}[H]\centering\small\rowcolors{2}{TableauPair}{TableauImpair}
\begin{tabular}{p{4.5cm}p{3.5cm}p{6.5cm}}
\toprule
\tableauHeader{Terme FR} & \tableauHeader{EN} & \tableauHeader{Définition} \\
\midrule
HIDS & Host-based IDS &
Système de détection d'intrusion basé sur l'hôte~; analyse les activités
locales (logs, fichiers, processus). \\
Horodatage & Timestamp &
Enregistrement de la date et heure d'un événement~; fondamental pour
la reconstruction de timeline. \\
Image disque & Disk image &
Copie bit à bit d'un support de stockage~; inclut les secteurs libres
et les données supprimées. \\
Incident de sécurité & Security incident &
Événement ou série d'événements confirmés portant atteinte à la CIA
(\textit{Confidentiality, Integrity, Availability}) du SI. \\
Indicateur de compromission & Indicator of compromise (IoC) &
Artefact observable indiquant une compromission passée ou en cours~:
hash, IP, domaine, chemin de fichier. \\
Injection de processus & Process injection &
Insertion de code malveillant dans l'espace d'adressage d'un processus
légitime pour masquer son exécution. \\
\bottomrule
\end{tabular}
\end{table}

% ─── L-M ──────────────────────────────────────────────────────────────────────
\section*{L -- M}

\begin{table}[H]\centering\small\rowcolors{2}{TableauPair}{TableauImpair}
\begin{tabular}{p{4.5cm}p{3.5cm}p{6.5cm}}
\toprule
\tableauHeader{Terme FR} & \tableauHeader{EN} & \tableauHeader{Définition} \\
\midrule
\textit{Living off the Land} & LotL / LOLBins &
Utilisation d'outils légitimes du système (powershell, certutil, wmic)
à des fins malveillantes pour contourner les outils de détection. \\
MACE & MACE (timestamps) &
Acronyme des quatre horodatages NTFS~: Modified, Accessed, Created,
Entry Modified. \\
Malware sans fichier & Fileless malware &
Code malveillant s'exécutant entièrement en mémoire, sans persistance
sur disque~; invisible aux antivirus basés sur les signatures de fichiers. \\
MFT & Master File Table &
Structure centrale de NTFS contenant un enregistrement pour chaque fichier
et répertoire du volume. \\
MISP & Malware Information Sharing Platform &
Plateforme open-source de partage et de gestion des indicateurs
de compromission. \\
MITRE ATT\&CK & MITRE ATT\&CK &
Base de connaissance mondiale des TTPs d'adversaires réels,
structurée en 14~tactiques et $>$200~techniques. \\
MTTD & Mean Time To Detect &
Durée moyenne entre la survenance d'un incident et sa détection. \\
MTTR & Mean Time To Respond &
Durée moyenne entre la détection d'un incident et les premières mesures
de confinement. \\
Mouvement latéral & Lateral movement &
Déplacement de l'attaquant vers d'autres systèmes dans le réseau compromis
après un accès initial. \\
\bottomrule
\end{tabular}
\end{table}

% ─── N-P ──────────────────────────────────────────────────────────────────────
\section*{N -- P}

\begin{table}[H]\centering\small\rowcolors{2}{TableauPair}{TableauImpair}
\begin{tabular}{p{4.5cm}p{3.5cm}p{6.5cm}}
\toprule
\tableauHeader{Terme FR} & \tableauHeader{EN} & \tableauHeader{Définition} \\
\midrule
Normalisation (logs) & Log normalization &
Conversion des journaux hétérogènes dans un format commun (ECS, CEF)
pour permettre la corrélation. \\
Ordre de volatilité & Order of volatility &
Hiérarchie de collecte des preuves en fonction de leur durée de vie~:
RAM~$\to$~réseau~$\to$~disque~$\to$~archives. \\
\textit{Parsing} & Parsing &
Décomposition d'une ligne de log en champs structurés. \\
Persistance & Persistence &
Mécanisme permettant à un malware de survivre aux redémarrages~:
Run keys, services, cron, authorized\_keys. \\
Playbook & Playbook &
Document décrivant la procédure logique de réponse à un type d'incident.
Répond à~: \og{}Quoi faire~?\fg{} \\
Preuve numérique & Digital evidence &
Donnée électronique pouvant être présentée devant une juridiction comme
élément de preuve. \\
\textit{Process hollowing} & Process hollowing &
Technique d'évasion~: processus légitime démarré en mode suspendu~;
son code est remplacé par du code malveillant avant reprise. \\
\bottomrule
\end{tabular}
\end{table}

% ─── R-S ──────────────────────────────────────────────────────────────────────
\section*{R -- S}

\begin{table}[H]\centering\small\rowcolors{2}{TableauPair}{TableauImpair}
\begin{tabular}{p{4.5cm}p{3.5cm}p{6.5cm}}
\toprule
\tableauHeader{Terme FR} & \tableauHeader{EN} & \tableauHeader{Définition} \\
\midrule
Rançongiciel & Ransomware &
Logiciel malveillant chiffrant les données de la victime et exigeant
une rançon pour la clé de déchiffrement. \\
Rétablissement & Recovery &
Phase de réponse aux incidents consistant à restaurer les systèmes à
leur état opérationnel normal. \\
RETEX & Post-incident review &
Retour d'expérience structuré après un incident pour identifier les causes
profondes et améliorer le dispositif. \\
Règle de corrélation & Correlation rule &
Règle SIEM définissant la condition logique, la fenêtre temporelle et
le seuil déclenchant une alerte. \\
Runbook & Runbook &
Document technique détaillant les commandes exactes à exécuter lors d'une
réponse à incident. Répond à~: \og{}Comment faire~?\fg{} \\
SIEM & Security Information and Event Management &
Plateforme centralisant la collecte, la normalisation, la corrélation
et l'alerting des événements de sécurité. \\
\textit{Slack space} & Slack space &
Espace résiduel entre la fin d'un fichier et la fin de son cluster~;
peut contenir des fragments de données antérieures. \\
SOC & Security Operations Center &
Structure organisationnelle dédiée à la surveillance continue du SI,
à la détection des menaces et à la coordination des réponses aux incidents. \\
STIX & Structured Threat Information eXpression &
Format JSON standardisé (OASIS) pour représenter des objets de
threat intelligence. \\
\bottomrule
\end{tabular}
\end{table}

% ─── T-Z ──────────────────────────────────────────────────────────────────────
\section*{T -- Z}

\begin{table}[H]\centering\small\rowcolors{2}{TableauPair}{TableauImpair}
\begin{tabular}{p{4.5cm}p{3.5cm}p{6.5cm}}
\toprule
\tableauHeader{Terme FR} & \tableauHeader{EN} & \tableauHeader{Définition} \\
\midrule
Tactique (ATT\&CK) & Tactic &
Objectif tactique de l'adversaire~: \textit{pourquoi} il réalise une action
(ex.~: Persistence, Lateral Movement). \\
TAXII & Trusted Automated eXchange of Intelligence Information &
Protocole HTTP(S) pour la transmission automatique d'objets STIX. \\
Technique (ATT\&CK) & Technique &
\textit{Comment} l'adversaire atteint son objectif tactique
(ex.~: T1110.001~— Password Guessing). \\
Temps de séjour & Dwell time &
Voir \textit{Dwell time}. \\
\textit{Threat hunting} & Threat hunting &
Processus proactif de recherche de signaux de compromission latents,
conduit par hypothèses CTI. \\
Timeline & Timeline &
Reconstruction chronologique des événements d'un incident, issue de la
corrélation de sources hétérogènes. \\
TLP & Traffic Light Protocol &
Système de classification du partage CTI~: RED, AMBER, GREEN, CLEAR. \\
\textit{Timestomping} & Timestomping &
Modification frauduleuse des timestamps MACE d'un fichier pour
brouiller la timeline forensique. \\
TTP & Tactics, Techniques and Procedures &
Comportements, méthodes et procédures caractéristiques d'un acteur
de la menace. Niveau le plus durable pour la défense. \\
Triage & Triage &
Qualification d'une alerte~: vrai ou faux positif~? gravité~? urgence~? \\
Volatilité (données) & Data volatility &
Durée de vie des données numériques~: les données les plus éphémères
(RAM) doivent être collectées en priorité. \\
XDR & Extended Detection and Response &
Corrélation unifiée HIDS + NIDS + SIEM + CTI avec réponse automatisée. \\
\textit{Zero-day} & Zero-day &
Vulnérabilité inconnue du vendeur et sans correctif disponible au moment
de son exploitation. \\
\bottomrule
\end{tabular}
\end{table}
